\section{Markov property for simple SCMs}

By making use of the `acyclification', we can extend the global Markov property for IL-CBNs to a global Markov property for simple SCMs.
The difference is that the Markov property for SCMs is formulated in terms of $\sigma$-separation rather than $d$-separation.
With the help of this Markov property, we can derive a very analogous theory for simple SCMs as for IL-CBNs, with a do-calculus and adjustment.

\subsection{Acyclifications}

In Section~\ref{sec:graph_acyclifications}, we defined the acyclification of a CDMG.
We can also define an operation with the same name on SCMs.
\begin{Def}[Acyclification of SCM]\label{def:scm_acyclification}
  Let $M = \lt \Sin, \Send, \Srnd, \CX, P, f \rt$ be a simple SCM with graph $G = G(M)$.
  For each strongly connected component $C$ of $G(M)$ (i.e., a set of the form
  $\Sc^G(v)$ for $v \in \Send$), let 
  $g_C : \CX_{\Sin \cup (\Send \setminus C) \cup \Srnd} \to \CX_C$
  be the unique solution function for $M_{\doit(\Send \setminus C)}$.
  Define $\tilde{f} : \CXJ \times \CXV \times \CXW \to \CXV$ by its components
  $$\tilde{f}_v(\xJ,\xV,\xW) = (g_{\Sc^G(v)})_v(\xJ,x_{\Send\setminus C},\xW)$$
  for $v \in V$.
  $M^\acy = \ltuple \Sin, \Send, \Srnd, \CX, P, \tilde{f} \rtuple$ is called the
  \emph{acyclification of $M$}.
\end{Def}
The crucial property of this definition is the following result, which also motivates its name.
\begin{Prp}\label{prp:scm_acyclification_properties}
  Let $M = \lt \Sin, \Send, \Srnd, \CX, P, f \rt$ be a simple SCM. Its acyclification $M^\acy$ is
  acyclic and observationally equivalent to $M$.
\end{Prp}
\begin{proof}
  We construct a directed graph $S$ from $G$ with its strongly connected components $\{\Sc^G(v) : v \in \Send \cup \Sin \cup \Srnd\}$
  as nodes, and directed edges
  $C \tuh D$ if there is a directed edge $c \tuh d$ in $G$ with $c \in C, d \in D$ and $C \ne D$.
  The graph $S$ cannot contain a directed cycle, as that would imply the existence of a directed
  cycle in $G$ that traverses more than one of its strongly connected components. Hence $S$ is a DAG.
  
  Choose a topological ordering $<$ of $S$.
  Any node $C$ in $S$ can only have incoming directed edges
  in $S$ from $\Pred^S_<(C)$. This implies that for $v \in V$, $C = \Sc^G(v)$ can only have incoming
  edges in $G$ from $\bigcup \Pred^S_\le(C)$. That implies that the causal mechanism $f_C$
  can only depend on variables in $\bigcup \Pred^S_\le(C)$, and hence
  the unique solution function $g_C$, and therefore $\tilde{f}_C$, can depend on variables in $\bigcup \Pred^S_<(C)$ only.
  Therefore, for $v,w \in V$, a directed edge $w \tuh v$ in $G(M^\acy)$ implies 
  $w \in \bigcup \Pred^S_<(\Sc^G(v))$. We can therefore refine the topological ordering $<$ of $S$ 
  to a topological ordering of $G(M^\acy)$, by arbitrarily ordering the nodes within each strongly connected component
  of $G$. Hence $G(M^\acy)$ is acyclic.

  $M$ and $M^\acy$ are observationally equivalent by construction: for all $x \in \CX$,
  \begin{align*}
  & x = \tilde{f}(x) \\
  & \iff \forall C \in S \cap V: x_C = \tilde{f}_C(x) \\
  & \iff \forall C \in S \cap V: x_C = g_C(x_J,x_{V\setminus C},x_W) \\
  & \iff \forall C \in S \cap V: x_C = f_C(x) \\
  & \iff x = f(x).
  \end{align*}
\end{proof}
The notion of acyclification of an SCM is compatible with that of a graph:
\begin{Prp}\label{prp:acyclifications_compatible}
  Let $M$ be a simple SCM.
  Then $G(M^\acy) \subseteq G'$ for any acyclification $G'$ of $G(M)$.
\end{Prp}
\begin{proof}
  By definition, the two graphs have the same nodes (input nodes $J$ and output nodes
  $V \cup W$). $G(M^\acy)$ has no bidirected edges, but $G'$ might. If there is a 
  directed edge $i \tuh j$ in $G(M^\acy)$ with $i \in J \cup V \cup W$ and $j \in V$, then
  the solution function $g_{\Sc^G(j)}$ of $M_{\doit(V\setminus \Sc^G(j))}$ depends on $x_i$. 
  This can only happen if
  $i \notin \Sc^G(j)$ and $i$ is a parent of some $k$ according to $M$ with $k \in \Sc^G(j)$,
  i.e., if $i \tuh k$ in $G(M)$. In that case, $i \tuh j$ in $G'$ by definition of
  the graphical acyclification.
\end{proof}
Hence, two nodes in the same strongly connected component of $G(M)$ do not have any edge between them
in $G(M^\acy)$, whereas they necessarily have a connecting edge in any acyclification $G'$ of $G(M)$.
For two nodes in different strongly connected components of $G(M)$, the edges in $G(M^\acy)$ are also
present in $G'$, but not necessarily vice versa, as some parent-relations may cancel out in the solution function.

\subsection{Global Markov property for simple SCMs}

With the help of the acyclifications,
we can easily derive a Markov property for simple SCMs from the Markov property for CBNs by reducing
the general cyclic case to an acyclic case.
\begin{Cor}[Global Markov property for simple SCMs]\label{cor:gmp_simple_scm}
  Let $M = \lt \Sin, \Send, \Srnd, \CX, P, f \rt$ be a simple SCM with graph $G(M)$
  and observational Markov kernel $P_{M}(X_V,X_W \given \doit(X_J))$.
  Then for all $A, B, C \ins J \cup V \cup W$ (not necessarily disjoint) we have the implication:
  $$A \Perp^\sigma_{G(M)} B \given C \implies X_A \Indep_{P_{M}(X_V,X_W \given \doit(\XJ))} X_B \given X_C.$$
  If one wants to make the implicit dependence on $J$ more explicit one can equivalently also write:
  $$A \Perp^\sigma_{G(M)} J \cup B \given C \implies X_A \Indep_{P_{M}(X_V,X_W \given \doit(\XJ))} X_J,X_B \given X_C.$$
\end{Cor}
\begin{proof}
  Choose an acyclification $G'$ of $G(M)$. Then:
  \begin{equation*}\begin{split}
    A \Perp^\sigma_{G(M)} B \given C 
    & \iff A \Perp^d_{G'} B \given C \\
    & \implies A \Perp^d_{G(M^\acy)} B \given C \\
    & \implies X_A \Indep_{P_{M^\acy}(X_V,X_W \given \doit(\XJ))} X_B \given X_C \\
    & \iff X_A \Indep_{P_{M}(X_V,X_W \given \doit(\XJ))} X_B \given X_C.
  \end{split}\end{equation*}
  For the various implications / equivalences, we used:
  \begin{enumerate}
    \item $G'$ is an acyclification of $G(M)$ together with Proposition~\ref{prp:sigma_sep_as_d_sep}; 
    \item $G(M^\acy) \subseteq G'$ from Proposition~\ref{prp:acyclifications_compatible}, and that removing edges cannot turn a $d$-separation into a $d$-connection; 
    \item the global Markov property Theorem~\ref{thm-gmp-mI-CBN} for $M^\acy$ interpreted as a causal Bayesian network (with deterministic Markov kernels for the endogenous variables, and purely probabilistic Markov kernels for the exogenous random variables), exploiting Proposition~\ref{prp:scm_acyclification_properties} that states that the acyclification $M^\acy$ is acyclic,
    \item by Proposition~\ref{prp:scm_acyclification_properties}, the acyclification $M^\acy$ has the same observational Markov kernel as the original SCM $M$.
  \end{enumerate}
\end{proof}
%This Markov property is not complete: it does not exploit all information in $G(M)$ that it possibly could to deduce 
%conditional independences. For example, endogenous variables without parents (for which the corresponding outcomes must be
%constant) could additionally be taken to block walks.
%\Joris{Should we change the condition for a non-collider to block the path if it is a deterministic function of $C$? We could detect this graphically by distinguishing deterministic kernels from non-deterministic ones (e.g.\ by a double circle).}

\subsection{Do-calculus for simple SCMs}

\Joris{TODO: update a la Patrick}

With the global Markov property for simple SCMs, it becomes straightforward to derive the do-calculus for simple SCMs.
First we will introduce some notation. The setting will be that a simple SCM $M = \lt \Sin, \Send, \Srnd, \CX, P, f \rt$
is given. 
As an auxiliary way to model hard interventions $\doit(B)$ for $B \subseteq V \cup W$ \Joris{if allowing intervention nodes
for exogenous input nodes, $B \subseteq V \cup W \cup J$, the following needs to be updated; perhaps also don't do this 
inline but have a proper definition} that leads to easy rules in the
do-calculus, we introduce additional intervention variables
$I_b$ for $b \in B$. We denote $I_B := (I_b)_{b \in B}$.
This gives an extended SCM $\tilde{M} = \lt \Sin \cup I_B, \Send, \Srnd, \CX \times \prod_{b \in B} (\CX_b \dot\cup \{\star\}), P, \tilde{f} \rt$ with causal mechanism with components
$$\tilde{f}_b(x_{V \cup J \cup W},x_{I_B}) := \begin{cases}
  f_b(x_{V \cup J \cup W}) & x_{I_b} = \star \\
  x_{I_b} & x_{I_b} \in \CX_b
\end{cases}$$
for $b \in B$, and $\tilde{f}_v(x) := f_v(x_{V \cup J \cup W})$ for $v \in V \setminus B$.
Here, $x_{I_b} = \star$ encodes that there is no intervention on $b$, 
while $x_{I_b} \ne \star$ encodes that the hard intervention $\doit(X_b = x_{I_b})$ is performed.
We will denote the graph of $M$ by $G$ and the graph of $\tilde{M}$ by $G_{\doit(I_B)}$.
In the do-calculus, we make use of the extended graph $G_{\doit(I_B)}$ to test the separation statement, while the conclusion
about the existence and identity of particular Markov kernels concerns those of the original SCM $M$.

Remember the notation of Markov kernels for simple SCMs from Definition~\ref{def:simple_scm_Markov_kernels}.
Its observational Markov kernel is denoted $P_{M}(X_{V\cup W} \mid \doit(X_J))$.
For a subset $T \subseteq V \cup W$, we write 
$$P_{M}(X_{(V\cup W)\setminus T} \mid \doit(X_J,X_T)) := P_{M_{\doit(T)}}(X_{(V\cup W)\setminus T} \mid \doit(X_J,X_T)).$$
By conditioning on a subset $S \subseteq (V \cup W) \setminus T$, we obtain the conditional Markov kernel
$$P_{M}(X_{(V\cup W)\setminus (T \cup S)} \mid \doit(X_J,X_T), X_S).$$

The only modification to the do-calculus for simple SCMs as compared to that for causal Bayesian networks is that
we have to replace all $d$-separations by $\sigma$-separations.
\begin{Cor}[Do-calculus for simple SCMs]\label{cor:scm_do_calculus}
  Let $M = \lt \Sin, \Send, \Srnd, \CX, P, f \rt$ be a simple SCM with graph $G = G(M)$.
  Let $A,B,C \ins V \cup W$ and $D \ins V \cup W \cup J$ be such that $A,B,C,D$ are pairwise disjoint.
  \Joris{In rule 3, should we allow for $B \subseteq V \cup W \cup J$?}
  Then we have the following 3 rules relating marginal conditional interventional Markov kernels of $M$:
    \begin{enumerate}
        \item \emph{Insertion/deletion of observation}: If we have:
            \[ A \Perp^\sigma_{G_{\doit(D)}} B \given C \cup D,   \]
            then there exists a Markov kernel:
            \[ P\lp X_A|\cancel{X_B},X_C,\doit(X_{D\cup \cancel{J}})\rp:\;  \Xcal_C \times \Xcal_{D} \dshto \Xcal_A \]
            that is a version of:
            \[P_{M}\lp X_A|X_{B_2},X_C,\doit(X_{D\cup J})\rp,\] 
            for every subset $B_2 \ins B$ simultaneously.
            In short we could write:
            \[ P_{M}\lp X_A|X_B,X_C,\doit(X_{D\cup J})\rp = P_{M}\lp X_A|X_C,\doit(X_D)\rp.\]
            %\[ P\lp X_A|X_B,X_C,\doit(X_D)\rp =P\lp X_A|X_{B_2},X_C,\doit(X_D)\rp = P\lp X_A|X_C,\doit(X_D)\rp.\]
            Note that this Markov kernel is only dependent on $x_C$ and $x_{D}$, and constant in $x_{J\setminus D}$.
            Such a Markov kernel is unique up to $P_{M}\lp X_B,X_C|\doit(X_{D\cup J})\rp$-null sets.
            %then we get the equality:
            %\[ P\lp X_A|X_B,X_C,\doit(X_D)\rp = P\lp X_A|X_C,\doit(X_D)\rp,  \]
            %$P\lp X_B,X_C|\doit(X_D)\rp$-a.s., to be read as that there exists a version of the rhs 
            %that is also a version of the lhs. 
         \item \emph{Action/observation exchange}: If we have:
            \[ A \Perp^\sigma_{G_{\doit(I_B,D)}} I_B \given B \cup C \cup D,   \]
            then there exists a Markov kernel:
            \[P\lp X_A|\cancel{\doit}(X_B),X_C,\doit(X_{D\cup \cancel{J}})\rp: \;\Xcal_B \times \Xcal_C \times \Xcal_{D} \dshto \Xcal_A\]
            that is a version of:            
            \[ P_{M}\lp X_A|X_{B_1},\doit(X_{B_2}),X_C,\doit(X_{D\cup J})\rp, \]
            for every decomposition $B=B_1\dcup B_2$ simultaneously. In short we could write:
%            \[ P_{M}\lp X_A|\doit(X_B),X_C,\doit(X_{D\cup J})\rp  =  P_{M}\lp X_A|X_B,X_C,\doit(X_{D\cup J})\rp.\]
%            or even
            \[ P_{M}\lp X_A|\doit(X_B),X_C,\doit(X_D)\rp  =  P_{M}\lp X_A|X_B,X_C,\doit(X_D)\rp.\]
            since these Markov kernels are constant in $x_{J\setminus D}$.
            Such a Markov kernel is unique up to $P_{M}\lp X_B,X_C|\doit(X_{I_B},X_{D\cup J})\rp$-null sets.
         \item \emph{Insertion/deletion of action}: If we have:
            \[ A \Perp^\sigma_{G_{\doit(I_B,D)}} I_B \given C  \cup D,   \]
            then there exists a Markov kernel: 
            \[P\lp X_A|\cancel{\doit(X_B)},X_C,\doit(X_{D\cup \cancel{J}})\rp:\; \Xcal_C \times \Xcal_{D} \dshto \Xcal_A\] 
            that is a version of:  
            \[ P_{M}\lp X_A|\doit(X_{B_2}),X_C,\doit(X_{D\cup J})\rp\]
            for every subset $B_2 \ins B$ simultaneously. In short we could write:
            \[ P_{M}\lp X_A|\doit(X_B),X_C,\doit(X_{D\cup J})\rp = P_{M}\lp X_A|X_C,\doit(X_{D})\rp.\]
            Note that this Markov kernel is only dependent on $x_C$ and $x_{D}$, and constant in $x_{J\setminus D}$.
            Such a Markov kernel is unique up to $P_{M}\lp X_C|\doit(X_{I_B},X_{D\cup J})\rp$-null set.
%          \item \Joris{\emph{Deletion of action II}: Here we take $B \subseteq J$ disjoint from $A,C,D$.
%            If we have:
%            \[ A \Perp^\sigma_{G_{\doit(D)}} \emptyset \given C \cup D, \]
%            then there exists a Markov kernel: 
%            \[P\lp X_A|X_C,\doit(X_{D\cup \cancel{J}})\rp:\; \Xcal_C \times \Xcal_{D} \dshto \Xcal_A\] 
%            that is a version of:  
%            \[ P_{M}\lp X_A|X_C,\doit(X_{D\cup J})\rp\]
%            In short we could write:
%            \[ P_{M}\lp X_A|X_C,\doit(X_{D\cup J})\rp = P_{M}\lp X_A|X_C,\doit(X_{D})\rp.\]
%            Note that this Markov kernel is only dependent on $x_C$ and $x_{D}$, but constant in $x_{J\setminus D}$.
%            Such a Markov kernel is unique up to $P_{M}\lp X_C|\doit(X_{D\cup J})\rp$-null set.}
    \end{enumerate}
\end{Cor}
\begin{proof}
  The proof is also analogous to that of Corollary~\ref{cor:do-calculus}, except that it applies the global
  Markov property for simple SCMs, Corollary~\ref{cor:gmp_simple_scm}, instead of the one for causal Bayesian networks,
  Theorem~\ref{thm-gmp-mI-CBN}. 
  %Furthermore, we do not assume that $J \subseteq D$, so therefore in each Markov kernel
  %from the SCM (but not the ones whose existence is part of the conditional independence statement) we replaced 
  %$\doit(X_D)$ by $\doit(X_{D \cup J})$.

%\Joris{Rule 4. The assumption:
%            \[ A \Perp^\sigma_{G_{\doit(D)}} B \given C  \cup D,   \]
%implies the conditional independence by GMP \ref{cor:gmp_simple_scm}:
%  \[ X_A \Indep_{P_{M}(X_V,X_W|\doit(X_{D\cup J}))} X_B \given X_C,X_D.\]
%So we have the following factorization:
%  \[ P_{M}\lp X_A,X_C|\doit(X_{D\cup J})\rp = Q(X_A|X_C,X_D) \otimes P_{M}\lp X_C|\doit(X_{D\cup J})\rp,\]
%for some Markov kernel $Q(X_A|X_C,X_D)$, which serves as a version of the conditional Markov kernel:
%  \[P_{M}\lp X_A|X_C,\doit(X_{D\cup J})\rp,\]
%and which is constant in $X_{B}$ and in $X_{J\setminus D}$.
%So we could write
%  \[P_{M}\lp X_A|X_C,\doit(X_{D\cup J})\rp = P_{M}\lp X_A|X_C,\doit(X_D)\rp\]
%}
\end{proof}

Thus, by making use of the do-calculus, 
we can in some cases derive the existence of Markov kernels $P_{M}(X_A \given \doit(X_B))$
where $J \not\subseteq B$;
these should be interpreted as Markov kernels $\CX_{J \cup B} \dshto \CX_{A}$ that are constant in 
$x_{J\setminus B}$. Essentially, the do-calculus follows immediately from application of the global Markov property. 
We do not know if this do-calculus is complete in our setting: in some cases, the global Markov property could perhaps
be used directly to obtain stronger conclusions.

\begin{Rem}
  In case you are wondering how to use rule 3 to insert or delete an action for an exogenous input
  variable $j \in J$: this can be done via the special case $B = I_B = \emptyset$ and $D = J \setminus \{j\}$.
  \Joris{This would become redundant}
\end{Rem}

%\Joris{
%\begin{Rem}
%I conjecture, for $B \subseteq V \cup W \cup J$ disjoint from $A \cup C \cup D$:
%\[ A \Perp^\sigma_{G_{\doit(I_{B\setminus J},D)}} J \given C \cup D \iff
%\forall B_2 \subseteq B \setminus J: A \Perp^\sigma_{G_{\doit(B_2,D)}} B \given C \cup D, \]
%\end{Rem}
%}

\begin{comment}
  We make use of the global Markov property (GMP) for simple SCMs, Corollary~\ref{cor:gmp_simple_scm}.\\

Point 1.) The assumption:
            \[ A \Perp^\sigma_{G_{\doit(D)}} B \given C \cup D,   \]
implies the conditional independence by GMP \ref{cor:gmp_simple_scm}:
  \[ X_A \Indep_{P(X_V,X_W|\doit(X_{D\cup J}))} X_B \given X_C,X_D.\]
So  we get the following factorization, where we can omit the deterministic variables from $X_D$ on the left of the conditioning lines:
  \[ P\lp X_A,X_B,X_C|\doit(X_{D\cup J})\rp = Q(X_A|X_C,X_D) \otimes P\lp X_B,X_C|\doit(X_{D\cup J})\rp,\]
for some Markov kernel $Q(X_A|X_C,X_D)$. Here
$Q(X_A|X_C,X_D)$ serves as a version of the conditional Markov kernel:
  \[P\lp X_A|X_B,X_C,\doit(X_{D\cup J})\rp.\]
If we marginalize out $X_{B_1}$ for any decomposition $B=B_1 \dcup B_2$ 
in the above factorization we also get:
  \[ P\lp X_A,X_C,X_{B_2}|\doit(X_{D\cup J})\rp = Q(X_A|X_C,X_D) \otimes P\lp X_C,X_{B_2}|\doit(X_{D\cup J})\rp,\]
showing that $Q(X_A|X_C,X_D)$ is also a version of:
  \[P\lp X_A|X_{B_2},X_C,\doit(X_{D\cup J})\rp.\]
In particular, this holds for $B_2 = \emptyset$.\\

Point 2.) The assumption:
            \[ A \Perp^\sigma_{G_{\doit(I_B,D)}} I_B \given B \cup C \cup D,   \]
implies the conditional independence by GMP \ref{cor:gmp_simple_scm}:
  \[ X_A \Indep_{P(X_V,X_W|\doit(X_{I_B},X_{D\cup J}))} X_{I_B} \given X_B,X_C,X_D.\]
So we have the following factorization:
  \[ P\lp X_A,X_B,X_C|\doit(X_{I_B},X_{D\cup J})\rp = Q(X_A|X_B,X_C,X_D) \otimes P\lp X_B,X_C|\doit(X_{I_B},X_{D\cup J})\rp,\]
for some Markov kernel $Q(X_A|X_B,X_C,X_D)$, which serves as a version of the conditional Markov kernel:
  \[P\lp X_A|X_B,X_C,\doit(X_{I_B},X_{D\cup J})\rp,\]
and which is independent of $X_{I_B}$.\\
We can now look at the different input values for any decomposition: $B=B_1 \dcup B_2$.
For this we put: $X_{I_{B_1}}= \star = (\star)_{v \in B_1}$ and  $ X_{I_{B_2}}=x_{B_2} \in \Xcal_{B_2}$.
This implies: 
\begin{align*} 
  & P\lp X_A,X_B,X_C|\doit(X_{{B_2}}=x_{B_2},X_{D\cup J})\rp \\
  &= P\lp X_A,X_B,X_C|\doit(X_{I_{B_1}}=\star,X_{I_{B_2}}=x_{B_2},X_{D\cup J})\rp \\
  &= Q(X_A|X_B,X_C,X_D) \otimes P\lp X_B,X_C|\doit(X_{I_{B_1}}=\star,X_{I_{B_2}}=x_{B_2},X_{D\cup J})\rp,\\
  &= Q(X_A|X_B,X_C,X_D) \otimes P\lp X_B,X_C|\doit(X_{{B_2}}=x_{B_2},X_{D\cup J})\rp.
\end{align*}
So we get:
  \[P\lp X_A,X_B,X_C|\doit(X_{B_2},X_{D\cup J})\rp 
  = Q(X_A|X_B,X_C,X_D) \otimes P\lp X_B,X_C|\doit(X_{B_2},X_{D\cup J})\rp,\]
where we can marginalize out the deterministic $X_{B_2}$:
  \[P\lp X_A,X_{B_1},X_C|\doit(X_{B_2},X_{D\cup J})\rp 
  = Q(X_A|X_B,X_C,X_D) \otimes P\lp X_{B_1},X_C|\doit(X_{B_2},X_{D\cup J})\rp.\]
This implies that $Q(X_A|X_B,X_C,X_D)$ is a version of the conditional Markov kernel:
  \[P\lp X_A|X_{B_1},X_C,\doit(X_{B_2},X_{D\cup J})\rp,\]
for every decomposition: $B=B_1 \dcup B_2$, simultaneously, in particular for the two cases
$B_2=\emptyset$ and $B_2 = B$.\\


Point 3.) The assumption:
            \[ A \Perp^\sigma_{G_{\doit(I_B,D)}} I_B \given C  \cup D,   \]
implies the conditional independence by GMP \ref{cor:gmp_simple_scm}:
  \[ X_A \Indep_{P(X_V,X_W|\doit(X_{I_B},X_{D\cup J}))} X_{I_B} \given X_C,X_D.\]
So we have the following factorization:
  \[ P\lp X_A,X_C|\doit(X_{I_B},X_{D\cup J})\rp = Q(X_A|X_C,X_D) \otimes P\lp X_C|\doit(X_{I_B},X_{D\cup J})\rp,\]
for some Markov kernel $Q(X_A|X_C,X_D)$, which serves as a version of the conditional Markov kernel:
  \[P\lp X_A|X_C,\doit(X_{I_B},X_{D\cup J})\rp,\]
and which is independent of $X_{I_B}$.\\ 
We can now look at the different input values for any decomposition: $B=B_1 \dcup B_2$.
For this we put: $X_{I_{B_1}}= \star = (\star)_{v \in B_1}$ and  $ X_{I_{B_2}}=x_{B_2} \in \Xcal_{B_2}$.
This implies:
\begin{align*} 
  & P\lp X_A,X_C|\doit(X_{{B_2}}=x_{B_2},X_{D\cup J})\rp \\
  &= P\lp X_A,X_C|\doit(X_{I_{B_1}}=\star,X_{I_{B_2}}=x_{B_2},X_{D\cup J})\rp \\
  &= Q(X_A|X_C,X_D) \otimes P\lp X_C|\doit(X_{I_{B_1}}=\star,X_{I_{B_2}}=x_{B_2},X_{D\cup J})\rp,\\
  &= Q(X_A|X_C,X_D) \otimes P\lp X_C|\doit(X_{{B_2}}=x_{B_2},X_{D\cup J})\rp.
\end{align*}
So we get:
\[P\lp X_A,X_C|\doit(X_{B_2},X_{D\cup J})\rp 
= Q(X_A|X_C,X_D) \otimes P\lp X_C|\doit(X_{B_2},X_{D\cup J})\rp,\]
which shows that $Q(X_A|X_C,X_D)$ is a version of the conditional Markov kernel:
\[P\lp X_A|X_C,\doit(X_{B_2},X_{D\cup J})\rp \]
for every decomposition: $B = B_1 \dcup B_2$, simultaneously, in particular for the two corner cases:
$B_2 = \emptyset$ and $B_2 = B$.
\end{proof}
\end{comment}

\subsection{Adjustment}

Let $M = \lt \Sin, \Send, \Srnd, \CX, P, f \rt$ be a simple SCM.
Let us assume $J \subseteq D$.
We are interested in estimating the conditional causal effect:
\[ P_{M}(X_A|X_C,\doit(X_B,X_D)), \]
but we only have data from:
\[ P_{M}(X_A,X_B,X_F|X_C,\doit(X_D)). \]
The following (pairwise disjoint) index sets will have the following roles:
\begin{description}
  \item[$A$]: the outcome variables of interest.
  \item[$B$]: the treatment or intervention variables.
  \item[$C$]: general conditional (context) variables under which the data was collected.
  \item[$D$]: general interventional (context) variables that were set by the experimenter, $J \ins D$.
  \item[$F_0$]: core adjustment variables, i.e.\ features that were measured.
  \item[$F_1$]: additional measured adjustment variables, with $F = F_0 \cup F_1$.
  \item[$H$]: additional unobserved variables.
\end{description}
\Joris{Add restrictions on what types of variables these are. E.g., $B \subseteq V \cup W$.}
We will make use of the same extended SCM $\tilde{M}$ with intervention variables $I_b$ for $b \in B$ and graph
$G_{\doit(I_B)}$ as for stating the do-calculus. \Joris{Better to refer to a formal than inline definition}

\begin{Ass}\label{ass:positivity}
  Let $M = \lt \Sin, \Send, \Srnd, \CX, P, f \rt$ be a simple SCM.
  Assume that:
  \begin{enumerate}
    \item Each measurable space $\Xcal_u$ for $u \in \Send \cup \Srnd$ comes equipped with a fixed measure $\mu_u$, and
    \item For every subset $D \ins \Send$ the interventional Markov kernel $P_M(X_D \mid \doit(X_{J \cup V \cup W \sm D}))$ is
      absolute continuous w.r.t.\ the product measure $\mu_D:=\bigotimes_{v \in D} \mu_v$ and vice versa:
            \[ \mu_D \ll P_M(X_D \mid \doit(X_{J \cup V \cup W \sm D})) \ll \mu_D. \]
  \end{enumerate}
%  The second point is satisfied if each interventional Markov kernel $P_M(X_u \mid \doit(X_{J \cup V \sm \{u\}}))$ has a strictly positive Doob-Radon-Nikodym derivative/density w.r.t.\ the measure $\mu_u$ for all $u \in \Send \cup \Srnd$.
%  \Joris{Is it? May not be due to cycles}
\end{Ass}
\begin{Thm}[General adjustment formula for simple SCMs]\label{thm:gen-adjustment-scms}
Given a simple SCM $M = \lt \Sin, \Send, \Srnd, \CX, P, f \rt$ with graph $G = G(M)$ that satisfies Assumption~\ref{ass:positivity},
assume that all the following conditions hold in the graph $G_{\doit(I_B,D)}$: % $G^+_{\doit(D)}$:
    \begin{align}
        (F_0 \cup H) & \Perp^\sigma_{G_{\doit(I_B,D)}} I_B \given (C \cup D), \\
%        H & \Perp^\sigma_{G_{\doit(I_B,D)}} I_B \given (C \cup D), \label{eq:adj-2} \\
        A & \Perp^\sigma_{G_{\doit(I_B,D)}} (F_1 \cup I_B) \given (B \cup F_0 \cup H \cup C \cup D), \\
%        A & \Perp^\sigma_{G_{\doit(I_B,D)}} F_1 \given ( B \cup F_0 \cup H \cup C \cup D),\label{eq:adj-4} \\
                      H &\Perp^\sigma_{G_{\doit(I_B,D)}} B \given (F \cup C \cup  I_B \cup D).
    \end{align}
Then we have the adjustment formula:
\[ P_{M}(X_A|X_C,\doit(X_B,X_D)) = 
    P_{M}(X_A|X_B,X_C,X_F,\doit(X_D)) \circ P_{M}(X_F|X_C,\doit(X_D)) \qquad \text{$\mu_{V\cup W}$-a.s.}
\]
\end{Thm}
\begin{proof}
  Analogous to that of Theorem~\ref{thm-gen-adjustment}, but now using
  the global Markov property for simple SCMs, Corollary~\ref{cor:gmp_simple_scm}, 
  instead of the one for causal Bayesian networks, Theorem~\ref{thm-gmp-mI-CBN}. 
\end{proof}
%Note that the ``a.s.'' qualifier is imprecise, a more precise statement could be obtained by tracing the uniqueness of the various Markov kernels appearing in the proof.

The following is just the special case $F_1=H=\emptyset$.
\begin{Cor}[Conditional backdoor covariate adjustment formula for simple SCMs]\label{thm:backdoor-cond-int-scms}
  Let $M = \lt \Sin, \Send, \Srnd, \CX, P, f \rt$ be a simple SCM with graph $G = G(M)$ that satisfies Assumption~\ref{ass:positivity}.
  Assume that the \emph{conditional backdoor criterion} in the graph $G_{\doit(I_B,D)}$ holds:
\begin{enumerate}
    \item $\displaystyle F \Perp^\sigma_{G_{\doit(I_B,D)}} I_B  \given (C \cup D)$, and:
    \item $\displaystyle A \Perp^\sigma_{G_{\doit(I_B,D)}} I_B  \given (B \cup F \cup C \cup D)$.
\end{enumerate}
Then we have the adjustment formula:
\[ P_{M}(X_A|X_C,\doit(X_B,X_D)) = 
    P_{M}(X_A|X_B,X_C,X_F,\doit(X_D)) \circ P_{M}(X_F|X_C,\doit(X_D)) \qquad \text{$\mu_{V\cup W}$-a.s.}
\]
\end{Cor}
The following is just the even more special case $C=D=J= \emptyset$.
\begin{Cor}[Backdoor covariate adjustment for simple SCMs]\label{thm:backdoor-scms}
 Given a simple SCM $M = \lt \Sin, \Send, \Srnd, \CX, P, f \rt$ with graph $G = G(M)$ that satisfies Assumption~\ref{ass:positivity},
 assume that the \emph{backdoor criterion} holds:
\begin{enumerate}
    \item $\displaystyle F \Perp^\sigma_{G_{\doit(I_B)}} I_B$, and:
    \item $\displaystyle A \Perp^\sigma_{G_{\doit(I_B)}} I_B  \given (B \cup F)$.
\end{enumerate}
Then we have the adjustment formula:
\[ P_{M}(X_A|\doit(X_B)) = 
    P_{M}(X_A|X_B,X_F) \circ P_{M}(X_F) \qquad \text{$\mu_{V\cup W}$-a.s.}
\]
\end{Cor}
%\Joris{The positivity conditions are pretty strong. Cannot we throw them overboard by 
%showing that
%\[ P_{M}(X_A|\doit(X_B)) = 
%P_{M}(X_A|X_B,X_F) \circ P_{M}(X_F) \qquad \text{$P(X_F,X_B)$-a.s. ?}
%\]
%Patrick notes that the latter thing is not really formally defined for the Markov kernel on the l.h.s..
%I countered that perhaps we should then define this notion, perhaps just by extending the one on the l.h.s.\
%by making it constant in $X_F$, i.e., as $P_{M}(X_A|\doit(X_B,\cancel{X_F}))$.
%At closer inspection, it seems that this is not such a good idea at all.
%However, what we could do in practice (in the discrete case) is to propagate bounds for anything that cannot be estimated from the data.
%}
